% Generated draft: requires mathematical and source review.
\documentclass[11pt]{article}
\usepackage[margin=25mm]{geometry}
\usepackage{fontspec}
\setmainfont{Latin Modern Roman}
\usepackage{amsmath,amssymb,mathtools,unicode-math}
\setmathfont{Latin Modern Math}
\usepackage{enumitem,needspace,xurl,hyperref}
\setlength{\parindent}{0pt}
\setlength{\parskip}{5pt}
\setlength{\emergencystretch}{3em}
\newcommand{\N}{\mathbb{N}}
\newcommand{\Z}{\mathbb{Z}}
\newcommand{\Q}{\mathbb{Q}}
\newcommand{\R}{\mathbb{R}}
\newcommand{\C}{\mathbb{C}}
\newcommand{\F}{\mathbb{F}}
\newcommand{\A}{\mathbb{A}}
\newcommand{\PP}{\mathbb{P}}
\newcommand{\E}{\mathbb{E}}
\newcommand{\eps}{\varepsilon}
\newcommand{\dd}{\,\mathrm{d}}
\begin{document}
\section*{UnsolvedMath ID: 30000693}
\textbf{Source problem code:} catalogue entry
\subsection{Definitions and mathematical statement}
For the finite-point comparison, let \(d\geq2\), let \(Y\subset\R^d\) be finite with \(\operatorname{aff}(Y)=\R^d\), and let \(H\) be an affine hyperplane. The source calls \(H\) simple when there exist \(q\in H\cap Y\) and an affine \((d-2)\)-flat \(F\subset H\) such that
\[
 (H\cap Y)\setminus\{q\}\subseteq F.
\]
It calls \(H\) not too simple when \(\operatorname{aff}(H\cap Y)=H\). Together these force the exceptional point off the containing flat; they do not merely ask for a small intersection or require exactly \(d\) intersection points. The original passage reports that such an \(H\) exists for every finite spanning \(Y\). This finite-point result is context attributed to the source, not a new theorem proved here.

The requested extension is to a \(k\)-dimensional real algebraic variety \(X\subseteq\R^d\) whose real point set affinely spans \(\R^d\). It asks first for suitable analogues of simple and not-too-simple for \(X\cap H\), and then for conditions on \(\deg X,k,d\) guaranteeing an affine hyperplane section with both properties. The higher-dimensional predicates are themselves part of the problem. The finite all-but-one-point definition cannot be applied literally to a positive-dimensional intersection as a substitute for choosing the requested analogue.

To become a precise existence statement, one must specify what simplicity measures (for example an explicitly chosen algebraic, geometric or topological condition), which condition excludes trivial sections, and how the chosen notions relate to the finite case. No Betti-number cutoff, smoothness condition, component bound or particular section topology is supplied, so none is selected. The variety conventions, irreducibility/purity, real versus complex dimension and degree definition are also absent. The source's affine-spanning hypothesis is retained; it is not replaced by projective nondegeneracy without stating a passage to projective closure.

The verification field explicitly marks the reconstruction unverified and says the official report defines these words only for finite point sets, leaving the higher-dimensional analogue open. This prevents a unique more specific conjecture from being recovered without invention. The literature triage checked 2026-08-21 reports related point-set and generic-hyperplane results but no matching general theorem. Its open label remains dated metadata, not an independently checked current result or evidence for a particular unchosen analogue.
\subsection{Short English statement}
Generalize the source's simple spanning hyperplane for a finite set to positive-dimensional real algebraic varieties. Both the meaning of a simple nontrivial section and the degree/dimension hypotheses are to be determined. The source explicitly warns that no unique higher-dimensional predicate is yet specified.
\subsection{Sources}
\begin{itemize}
\item Local UnsolvedMath registry, permanent ID 30000693, source code catalogue entry, ``Simple Hyperplane Sections of Real Algebraic Varieties,'' including the finite definitions, original extension question and reconstructed-unverified qualification.
\item \emph{Geometric and Topological Combinatorics (2007)}, Oberwolfach Reports, pp.\ 195--272, DOI 10.4171/owr/2007/04: \url{https://doi.org/10.4171/owr/2007/04}.
\item Supplied reference ``higher-dimensional Sylvester problem'': \url{https://arxiv.org/abs/1608.03189}.
\item Supplied reference ``sections of varieties discussion'': \url{https://mathoverflow.net/questions/4661/the-sylvester-gallai-theorem-and-sections-of-varieties-with-simple-topology}.
\item Literature triage checked 2026-08-21; local record updated 2026-09-16. No external source has been visited or independently checked here.
\end{itemize}
\end{document}
